\chapter{第六章分离性}
\section{T1,T2,T0正则空间}
T1空间不一定是T2空间
\begin{theorem}[定理 6.2]
设 $X$ 为拓扑空间，则 $X$ 是 $T_1$ 空间当且仅当对任意 $x\in X$，单点集 $\{x\}$ 是 $X$ 中的闭集。
\end{theorem}
\begin{example}
设
\[
X=\{a,b,c\},\qquad 
\mathcal T=\{X,\varnothing,\{a\},\{b\},\{a,b\}\}.
\]
对于 $X$ 中的两点 $b$ 与 $c$，每个含 $c$ 的开集都含有 $a$ 与 $b$，
因此 $X$ 不是 $T_1$ 空间。
但对于点 $a$ 与 $b$，分别存在开集 $\{a\}$ 与 $\{b\}$，
使得 $c\notin\{a\}$ 且 $c\notin\{b\}$。
\end{example}
\begin{definition}[定义 6.4]
设 $X$ 为拓扑空间。若对 $X$ 中任意两个不同的点 $x,y$，
存在开集 $U_x$ 使得
\[
x\in U_x\subset X\setminus\{y\},
\]
或者存在开集 $U_y$ 使得
\[
y\in U_y\subset X\setminus\{x\},
\]
则称 $X$ 为 $T_0$ 空间。
\end{definition}

\begin{remark}
例 6.3 中的空间是 $T_0$ 空间，但不是 $T_1$ 空间。
\end{remark}
\begin{theorem}[都是T某空间的直积也是T某空间]
若对每个 $\alpha\in\Lambda$，空间 $X_\alpha$ 都是 $T_i$ 空间，
其中 $i\in\{0,1,2\}$，则其直积空间
\[
\prod_{\alpha\in\Lambda} X_\alpha
\]
也是 $T_i$ 空间，$i\in\{0,1,2\}$。
\end{theorem}
\begin{definition}[T1空间+只要x不属于闭集就能找到包含x的开集和包含F的开集不相交]
设 $X$ 为拓扑空间。若 $X$ 是 $T_1$ 空间，且对 $X$ 中任一闭集 $F$
及任意点 $x\notin F$，都存在 $X$ 中的开集 $U_x$ 与 $U_F$，使得
\[
x\in U_x,\qquad F\subset U_F,\qquad U_x\cap U_F=\varnothing,
\]
则称 $X$ 为\emph{正则空间}。
\end{definition}

\begin{theorem}[正则空间都是T2空间]
每个正则空间都是 $T_2$ 空间。
\end{theorem}

\begin{proof}
若 $X$ 是正则空间，则 $X$ 是 $T_1$ 空间。
对 $X$ 中任意两个不同的点 $x,y$，有 $x\notin\{y\}$ 且 $\{y\}$ 为闭集。
由正则性，存在 $X$ 中的开集 $U_x$ 与 $U_y$，使得
\[
x\in U_x,\qquad \{y\}\subset U_y,\qquad U_x\cap U_y=\varnothing.
\]
因此 $x\in U_x,\ y\in U_y$ 且 $U_x\cap U_y=\varnothing$，从而 $X$ 是 $T_2$ 空间。
\end{proof}

\begin{theorem}[T1空间+包含x的开集都能找到他的开子集包含x，开子集的闭包也在他里面
]
设 $X$ 为 $T_1$ 拓扑空间，则 $X$ 是正则空间当且仅当
对任意 $x\in X$ 及包含 $x$ 的任一开集 $U$，
都存在开集 $V_x$，使得
\[
x\in V_x\subset \overline{V_x}\subset U.
\]
\end{theorem}